\documentclass[11pt,a4paper]{article}

% ============================================================
% COURS COMPLET DE MÉCANIQUE DU SOLIDE
% Auteur : Pierre Garcia-Buscaylet
% ============================================================

\usepackage[utf8]{inputenc}
\usepackage[T1]{fontenc}
\usepackage[french]{babel}
\usepackage{lmodern}
\usepackage{microtype}

\usepackage{amsmath,amssymb,amsthm,mathtools}
\usepackage{siunitx}
\usepackage{tikz}
\usepackage{pgfplots}
\usepackage[most]{tcolorbox}
\usepackage{enumitem}
\usepackage{geometry}
\usepackage{fancyhdr}
\usepackage{hyperref}
\usepackage{xcolor}
\usepackage{array}
\usepackage{booktabs}
\usepackage{multicol}
\usepackage{caption}

\pgfplotsset{compat=1.18}
\usetikzlibrary{arrows.meta,calc,angles,quotes,patterns,decorations.pathmorphing}

% ============================================================
% MISE EN PAGE
% ============================================================

\geometry{margin=2.2cm}
\setlength{\parindent}{0pt}
\setlength{\parskip}{0.55em}

\hypersetup{
    colorlinks=true,
    linkcolor=blue!55!black,
    urlcolor=blue!55!black,
    citecolor=blue!55!black
}

\pagestyle{fancy}
\fancyhf{}
\lhead{\textsc{Mécanique du solide}}
\rhead{\textsc{Pierre Garcia-Buscaylet}}
\cfoot{\thepage}

\sisetup{
  locale = FR,
  per-mode=symbol,
  exponent-product=\cdot,
  output-decimal-marker={,}
}

% ============================================================
% COULEURS ET BOÎTES
% ============================================================

\definecolor{bleuprepa}{RGB}{18,55,115}
\definecolor{vertprepa}{RGB}{20,105,70}
\definecolor{rougeprepa}{RGB}{160,40,40}
\definecolor{orangeprepa}{RGB}{190,105,20}
\definecolor{violetprepa}{RGB}{95,55,130}
\definecolor{grisclair}{RGB}{245,247,250}

\tcbset{
  enhanced,
  breakable,
  sharp corners=downhill,
  boxrule=0.8pt,
  colback=white,
  fonttitle=\bfseries,
  before skip=0.8em,
  after skip=0.8em
}

\newtcbtheorem[number within=section]{definitionbox}{Définition}
{colframe=bleuprepa,colback=blue!3!white}{def}

\newtcbtheorem[number within=section]{propositionbox}{Proposition}
{colframe=vertprepa,colback=green!3!white}{prop}

\newtcbtheorem[number within=section]{theoremBox}{Théorème}
{colframe=rougeprepa,colback=red!3!white}{thm}

\newtcbtheorem[number within=section]{methodbox}{Méthode}
{colframe=orangeprepa,colback=orange!4!white}{meth}

\newtcbtheorem[number within=section]{examplebox}{Exemple}
{colframe=violetprepa,colback=violet!3!white}{ex}

\newtcbtheorem[number within=section]{warningbox}{Erreur classique}
{colframe=rougeprepa,colback=red!4!white}{warn}

\newtcbtheorem[number within=section]{exercisebox}{Exercice}
{colframe=black!70,colback=grisclair}{exo}

\newtcbtheorem[number within=section]{correctionbox}{Correction}
{colframe=vertprepa,colback=green!2!white}{corr}

\newtcbtheorem[number within=section]{intuitionbox}{Intuition physique}
{colframe=bleuprepa,colback=blue!2!white}{intu}

\newtcbtheorem[number within=section]{historicalbox}{Note historique}
{colframe=black!65,colback=yellow!8!white}{hist}

% ============================================================
% COMMANDES
% ============================================================

\newcommand{\R}{\mathbb{R}}
\newcommand{\vect}[1]{\overrightarrow{#1}}
\newcommand{\dd}{\,\mathrm{d}}
\newcommand{\der}[2]{\frac{\mathrm{d}#1}{\mathrm{d}#2}}
\newcommand{\dder}[2]{\frac{\mathrm{d}^2#1}{\mathrm{d}#2^2}}
\newcommand{\pder}[2]{\frac{\partial #1}{\partial #2}}
\newcommand{\norm}[1]{\left\lVert #1 \right\rVert}
\newcommand{\scal}[2]{#1\cdot #2}
\newcommand{\vprod}{\wedge}
\newcommand{\mat}[1]{\mathbf{#1}}
\newcommand{\base}{(\vec e_x,\vec e_y,\vec e_z)}
\newcommand{\Rep}{\mathcal R}
\newcommand{\Solide}{\mathcal S}
\newcommand{\niveau}[1]{\textbf{Niveau #1}}

% ============================================================
% DOCUMENT
% ============================================================

\begin{document}

\begin{titlepage}
\centering
\vspace*{1.5cm}

{\Huge\bfseries Cours complet de mécanique du solide\par}
\vspace{0.5cm}
{\LARGE Du point matériel au solide indéformable tridimensionnel\par}

\vspace{2cm}

\begin{tikzpicture}[scale=1.15]
  \draw[->,thick] (-0.5,0) -- (5,0) node[right] {$x$};
  \draw[->,thick] (0,-0.5) -- (0,3) node[above] {$z$};
  \draw[thick,fill=blue!8] (1.5,1.1) circle (0.75);
  \draw[thick,fill=orange!15] (3.6,1.1) rectangle (4.8,2.1);
  \draw[->,very thick,red] (1.5,1.1) -- (2.6,2.0) node[midway,above] {$\vec v$};
  \draw[->,very thick,vertprepa] (4.2,1.6) -- (4.2,2.8) node[midway,right] {$\vec \Omega$};
  \node at (1.5,0.1) {point matériel};
  \node at (4.2,0.1) {solide indéformable};
  \draw[->,thick,dashed] (2.25,1.1) -- (3.55,1.6);
  \node at (2.9,0.75) {discret $\to$ continu};
\end{tikzpicture}

\vfill

{\Large Auteur : Pierre Garcia-Buscaylet\par}
{\Large Date : \today\par}

\vfill

\begin{tcolorbox}[colframe=bleuprepa,colback=blue!3!white,title=Objectif du polycopié]
Ce document construit progressivement la mécanique du solide à partir de la mécanique du point :
point matériel, système discret, passage au continu, cinématique du solide indéformable,
géométrie des masses, tenseur d'inertie, dynamique, énergie, liaisons, roulement,
stabilité et problèmes contextualisés.
\end{tcolorbox}

\end{titlepage}

\tableofcontents
\newpage

% ============================================================
\section*{Sources publiques d'inspiration}
\addcontentsline{toc}{section}{Sources publiques d'inspiration}

Ce polycopié est une rédaction originale inspirée de thèmes classiques présents dans des ressources publiques :
cours de mécanique du solide de licence et d'école d'ingénieur, notes de CPGE sur les torseurs cinétiques et dynamiques, et notes universitaires sur le tenseur d'inertie et les équations d'Euler.

\begin{itemize}
  \item UGA / LPSC, \emph{Cinétique des solides indéformables}.
  \item Jean-Philippe Matas, Grenoble INP, \emph{Mécanique du solide}.
  \item MIT OpenCourseWare, \emph{Rigid Body Dynamics}, chapitres sur le tenseur d'inertie et les équations d'Euler.
  \item SILLAGES / Unisciel, \emph{Mécanique du solide}.
  \item Ressources de sciences industrielles CPGE sur la dynamique et l'énergétique des solides indéformables.
\end{itemize}

Les contenus ci-dessous ne recopient pas ces sources : ils réorganisent les notions et les démonstrations dans un cours cohérent, progressif et autonome.

\newpage

% ============================================================
\section{Notations, conventions et cadre}

\begin{definitionbox}{Référentiel}{}
Un référentiel $\Rep$ est la donnée d'un observateur, d'une horloge et d'un repère spatial.
On note souvent
\[
\Rep=(O,\vec e_x,\vec e_y,\vec e_z).
\]
Un référentiel est dit galiléen lorsqu'un point matériel isolé y possède un mouvement rectiligne uniforme.
\end{definitionbox}

\begin{definitionbox}{Vecteurs, moments et produits}{}
Dans tout le document :
\begin{itemize}
  \item les vecteurs sont notés avec une flèche : $\vec v$, $\vec F$, $\vec \Omega$ ;
  \item le produit scalaire est noté $\vec a\cdot\vec b$ ;
  \item le produit vectoriel est noté $\vec a\wedge\vec b$ ;
  \item le moment en $A$ d'une force $\vec F$ appliquée en $M$ vaut
  \[
  \vec M_A(\vec F)=\vect{AM}\wedge \vec F ;
  \]
  \item la dérivée temporelle dans un référentiel $\Rep$ est notée
  \[
  \left(\der{\vec u}{t}\right)_{\Rep}.
  \]
\end{itemize}
\end{definitionbox}

\begin{warningbox}{Confusion entre force et moment}{}
Une force tend à produire une translation. Un moment de force tend à produire une rotation.
Deux forces opposées peuvent avoir une résultante nulle mais un moment non nul : c'est un couple.
\end{warningbox}

\begin{definitionbox}{Solide indéformable}{}
Un solide $\Solide$ est dit indéformable si la distance entre deux points quelconques $A,B$ du solide reste constante :
\[
\forall A,B\in\Solide,\qquad AB=\text{constante}.
\]
Cela interdit toute déformation interne. Le mouvement d'un solide indéformable est donc entièrement décrit par la position d'un point et par l'orientation du solide.
\end{definitionbox}

\begin{intuitionbox}{Pourquoi passer du point au solide ?}{}
Un point matériel ne possède ni taille, ni forme, ni orientation. Il ne peut donc pas tourner sur lui-même.
Un solide réel, au contraire, possède une extension spatiale. Sa masse est répartie dans l'espace.
Cette répartition modifie sa résistance à la rotation : c'est le rôle du moment d'inertie puis du tenseur d'inertie.
\end{intuitionbox}

% ============================================================
\section{Chapitre 0 — Rappels de mécanique du lycée}

\subsection{Position, vitesse, accélération}

\begin{definitionbox}{Position, vitesse, accélération}{}
Pour un point matériel $M$ dans un référentiel $\Rep$, la position est donnée par le vecteur
\[
\vec r(t)=\vect{OM}(t).
\]
La vitesse et l'accélération sont
\[
\vec v(t)=\left(\der{\vec r}{t}\right)_{\Rep},
\qquad
\vec a(t)=\left(\der{\vec v}{t}\right)_{\Rep}
=\left(\dder{\vec r}{t}\right)_{\Rep}.
\]
\end{definitionbox}

\begin{center}
\begin{tikzpicture}[scale=1.1]
\draw[->,thick] (0,0)--(5,0) node[right]{$x$};
\draw[->,thick] (0,0)--(0,3) node[above]{$y$};
\coordinate (M) at (3.4,1.7);
\draw[->,very thick,blue] (0,0)--(M) node[midway,above]{$\vec r$};
\fill (M) circle (2pt) node[above right]{$M$};
\draw[->,very thick,red] (M)--++(1.0,0.6) node[right]{$\vec v$};
\draw[->,very thick,vertprepa] (M)--++(-0.4,0.9) node[above]{$\vec a$};
\end{tikzpicture}
\end{center}

\subsection{Lois de Newton}

\begin{theoremBox}{Principe fondamental de la dynamique du point}{}
Dans un référentiel galiléen, un point matériel de masse $m$ soumis à des forces $\vec F_1,\dots,\vec F_n$ vérifie
\[
m\vec a=\sum_{i=1}^{n}\vec F_i.
\]
\end{theoremBox}

\begin{methodbox}{Méthode de résolution d'un problème de mécanique du point}{}
\begin{enumerate}
  \item Définir le système étudié.
  \item Choisir un référentiel, idéalement galiléen.
  \item Choisir un repère adapté.
  \item Faire le bilan des forces.
  \item Projeter la relation $m\vec a=\sum\vec F$.
  \item Résoudre les équations différentielles.
  \item Interpréter les unités et l'ordre de grandeur.
\end{enumerate}
\end{methodbox}

\subsection{Forces usuelles}

\begin{center}
\begin{tabular}{lll}
\toprule
Force & Expression typique & Remarque \\
\midrule
Poids & $\vec P=m\vec g$ & Vertical vers le bas \\
Réaction normale & $\vec N$ & Perpendiculaire au support \\
Frottement solide & $\norm{\vec T}\leq \mu_s N$ & Statique ou dynamique \\
Ressort & $\vec F=-k(x-\ell_0)\vec e_x$ & Force de rappel \\
Traînée fluide linéaire & $\vec f=-\lambda \vec v$ & Régime de Stokes \\
Traînée quadratique & $\vec f=-\frac12\rho C_D S v\vec v$ & Grandes vitesses \\
\bottomrule
\end{tabular}
\end{center}

\subsection{Travail, puissance, énergie}

\begin{definitionbox}{Travail et puissance}{}
La puissance d'une force $\vec F$ appliquée à un point de vitesse $\vec v$ est
\[
\mathcal P=\vec F\cdot\vec v.
\]
Le travail entre deux instants est
\[
W=\int_{t_1}^{t_2}\mathcal P\,\dd t
=\int_{A}^{B}\vec F\cdot \dd \vec r.
\]
\end{definitionbox}

\begin{theoremBox}{Théorème de l'énergie cinétique}{}
Pour un point matériel de masse $m$ :
\[
\Delta E_c=\sum_i W(\vec F_i),
\qquad
E_c=\frac12 mv^2.
\]
\end{theoremBox}

\subsection{Mouvement circulaire}

Pour un mouvement circulaire de rayon $R$ et de vitesse angulaire $\omega$ :
\[
v=R\omega,
\qquad
a_n=\frac{v^2}{R}=R\omega^2.
\]

\begin{examplebox}{Ordre de grandeur}{}
Une roue de rayon $R=\SI{0.30}{m}$ tournant à $\omega=\SI{20}{rad.s^{-1}}$ possède une vitesse périphérique
\[
v=R\omega=\SI{6.0}{m.s^{-1}}.
\]
\end{examplebox}

\begin{exercisebox}{Applications directes — rappels}{}
\begin{enumerate}
  \item Une masse de \SI{2}{kg} est soumise à une force constante de \SI{6}{N}. Calculer son accélération.
  \item Une voiture prend un virage de rayon \SI{50}{m} à \SI{36}{km.h^{-1}}. Calculer son accélération normale.
  \item Un ressort de raideur \SI{100}{N.m^{-1}} est allongé de \SI{5}{cm}. Calculer la force exercée.
\end{enumerate}
\end{exercisebox}

% ============================================================
\section{Chapitre 1 — Mécanique du point matériel}

\subsection{Cinématique vectorielle}

En coordonnées cartésiennes :
\[
\vec r=x\vec e_x+y\vec e_y+z\vec e_z,
\qquad
\vec v=\dot x\vec e_x+\dot y\vec e_y+\dot z\vec e_z,
\qquad
\vec a=\ddot x\vec e_x+\ddot y\vec e_y+\ddot z\vec e_z.
\]

En coordonnées polaires planes :
\[
\vec r=r\vec e_r,
\qquad
\vec v=\dot r\vec e_r+r\dot\theta\vec e_\theta,
\]
\[
\vec a=(\ddot r-r\dot\theta^2)\vec e_r+(r\ddot\theta+2\dot r\dot\theta)\vec e_\theta.
\]

\begin{center}
\begin{tikzpicture}[scale=1.15]
\draw[->] (0,0)--(4,0) node[right]{$x$};
\draw[->] (0,0)--(0,3) node[above]{$y$};
\coordinate (M) at (3,2);
\draw[->,thick,blue] (0,0)--(M) node[midway,above]{$r$};
\draw[dashed] (3,0)--(M);
\draw pic[draw,->,"$\theta$",angle eccentricity=1.4] {angle=(4,0)--(0,0)--(M)};
\fill (M) circle (2pt) node[above]{$M$};
\draw[->,red,thick] (M)--++(0.83,0.55) node[right]{$\vec e_r$};
\draw[->,vertprepa,thick] (M)--++(-0.55,0.83) node[above]{$\vec e_\theta$};
\end{tikzpicture}
\end{center}

\subsection{Quantité de mouvement et moment cinétique}

\begin{definitionbox}{Quantité de mouvement}{}
La quantité de mouvement d'un point matériel est
\[
\vec p=m\vec v.
\]
La deuxième loi de Newton s'écrit alors
\[
\der{\vec p}{t}=\sum \vec F.
\]
\end{definitionbox}

\begin{definitionbox}{Moment cinétique d'un point}{}
Le moment cinétique en $O$ d'un point matériel $M$ est
\[
\vec L_O=\vect{OM}\wedge m\vec v.
\]
\end{definitionbox}

\begin{propositionbox}{Théorème du moment cinétique pour un point}{}
Dans un référentiel galiléen :
\[
\der{\vec L_O}{t}=\vec M_O^{\rm ext},
\qquad
\vec M_O^{\rm ext}=\vect{OM}\wedge \sum \vec F.
\]
\end{propositionbox}

\begin{proof}
On dérive :
\[
\der{}{t}(\vect{OM}\wedge m\vec v)
=
\vec v\wedge m\vec v+\vect{OM}\wedge m\vec a.
\]
Or $\vec v\wedge m\vec v=\vec 0$ et $m\vec a=\sum \vec F$.
Donc
\[
\der{\vec L_O}{t}=\vect{OM}\wedge\sum\vec F.
\]
\end{proof}

\begin{warningbox}{Moment par rapport à un point ou à un axe}{}
Le moment en un point est un vecteur :
\[
\vec M_O=\vect{OM}\wedge\vec F.
\]
Le moment par rapport à un axe orienté $\Delta$ de vecteur unitaire $\vec u$ est un scalaire :
\[
M_\Delta=\vec M_O\cdot\vec u,
\]
où $O$ est un point de l'axe.
\end{warningbox}

% ============================================================
\section{Chapitre 2 — Systèmes discrets de points matériels}

\subsection{Définitions générales}

On considère un système discret constitué de $N$ points matériels $M_i$ de masses $m_i$.

\begin{definitionbox}{Masse totale et centre de masse}{}
La masse totale vaut
\[
M=\sum_{i=1}^N m_i.
\]
Le centre de masse $G$ est défini par
\[
\vect{OG}=\frac{1}{M}\sum_{i=1}^{N}m_i\vect{OM_i}.
\]
\end{definitionbox}

\begin{center}
\begin{tikzpicture}
\foreach \x/\y/\m in {0/0/1,2/0.6/2,1.2/2/1.5,3/1.8/1}{
  \fill[blue] (\x,\y) circle (2pt);
  \node[above] at (\x,\y) {$m_{\m}$};
}
\fill[red] (1.55,1.05) circle (3pt) node[right]{$G$};
\draw[->,thick] (-0.5,-0.3)--(3.8,-0.3) node[right]{$x$};
\draw[->,thick] (-0.5,-0.3)--(-0.5,2.5) node[above]{$y$};
\end{tikzpicture}
\end{center}

\subsection{Quantité de mouvement totale}

\[
\vec P=\sum_{i=1}^{N}m_i\vec v_i.
\]

Comme
\[
M\vect{OG}=\sum_i m_i\vect{OM_i},
\]
en dérivant :
\[
M\vec v_G=\sum_i m_i\vec v_i=\vec P.
\]

\begin{theoremBox}{Théorème de la résultante cinétique}{}
Dans un référentiel galiléen :
\[
M\vec a_G=\sum \vec F_{\rm ext}.
\]
\end{theoremBox}

\begin{proof}
Pour chaque point :
\[
m_i\vec a_i=\vec F_i^{\rm ext}+\sum_{j\neq i}\vec F_{j\to i}^{\rm int}.
\]
En sommant :
\[
\sum_i m_i\vec a_i=\sum_i\vec F_i^{\rm ext}
+\sum_i\sum_{j\neq i}\vec F_{j\to i}^{\rm int}.
\]
Par action-réaction, les forces intérieures se compensent deux à deux. Donc
\[
M\vec a_G=\sum\vec F_{\rm ext}.
\]
\end{proof}

\subsection{Moment cinétique total}

\[
\vec L_O=\sum_i \vect{OM_i}\wedge m_i\vec v_i.
\]

\begin{theoremBox}{Théorème du moment cinétique pour un système}{}
Dans un référentiel galiléen :
\[
\left(\der{\vec L_O}{t}\right)_{\Rep}
=
\sum_i \vect{OM_i}\wedge\vec F_i^{\rm ext},
\]
si les forces intérieures sont centrales.
\end{theoremBox}

\subsection{Théorèmes de Koenig}

\begin{theoremBox}{Premier théorème de Koenig}{}
Le moment cinétique en $O$ se décompose en :
\[
\vec L_O=M\vect{OG}\wedge\vec v_G+\vec L_G^*,
\]
où $\vec L_G^*$ est le moment cinétique propre calculé dans le référentiel barycentrique.
\end{theoremBox}

\begin{theoremBox}{Second théorème de Koenig}{}
L'énergie cinétique totale se décompose en :
\[
E_c=\frac12 Mv_G^2+E_c^*.
\]
\end{theoremBox}

\begin{intuitionbox}{Sens physique}{}
Le mouvement total d'un système se décompose en deux parties :
\begin{itemize}
  \item le mouvement global de son centre de masse ;
  \item le mouvement interne autour du centre de masse.
\end{itemize}
Pour un solide indéformable, le mouvement interne est essentiellement une rotation.
\end{intuitionbox}

% ============================================================
\section{Chapitre 3 — Passage du discret au continu}

\subsection{De la somme à l'intégrale}

Pour un système discret :
\[
M=\sum_i m_i,
\qquad
\vect{OG}=\frac{1}{M}\sum_i m_i\vect{OM_i}.
\]

Pour un solide continu de masse volumique $\rho(M)$ :
\[
\dd m=\rho(M)\dd V,
\]
donc
\[
M=\iiint_{\Solide}\rho(M)\dd V,
\]
\[
\vect{OG}=\frac{1}{M}\iiint_{\Solide}\vect{OM}\rho(M)\dd V.
\]

\begin{definitionbox}{Passage discret-continu}{}
Le passage
\[
\sum_i m_i f(M_i)
\quad\longrightarrow\quad
\iiint_{\Solide}f(M)\rho(M)\dd V
\]
correspond à une limite de sommes de Riemann lorsque la taille des éléments de masse tend vers zéro.
\end{definitionbox}

\begin{examplebox}{Centre de masse d'une tige homogène}{}
Pour une tige homogène de longueur $L$, placée sur l'axe $Ox$ entre $0$ et $L$ :
\[
M=\int_0^L \lambda \dd x=\lambda L.
\]
Puis
\[
x_G=\frac{1}{M}\int_0^L x\lambda \dd x
=\frac{\lambda L^2/2}{\lambda L}
=\frac L2.
\]
\end{examplebox}

\subsection{Moments d'inertie usuels}

\begin{definitionbox}{Moment d'inertie par rapport à un axe}{}
Le moment d'inertie d'un solide $\Solide$ par rapport à un axe $\Delta$ est
\[
I_\Delta=\iiint_{\Solide} r_\perp^2\,\dd m,
\]
où $r_\perp$ est la distance du point matériel élémentaire à l'axe.
\end{definitionbox}

\begin{center}
\begin{tabular}{lll}
\toprule
Solide homogène & Axe & Moment d'inertie \\
\midrule
Tige longueur $L$ & centre, perpendiculaire & $\frac{1}{12}ML^2$ \\
Disque rayon $R$ & axe de symétrie & $\frac12 MR^2$ \\
Cylindre plein rayon $R$ & axe de symétrie & $\frac12 MR^2$ \\
Cylindre creux mince & axe de symétrie & $MR^2$ \\
Sphère pleine rayon $R$ & diamètre & $\frac25 MR^2$ \\
Parallélépipède $a,b,c$ & axe $x$ par centre & $\frac{1}{12}M(b^2+c^2)$ \\
\bottomrule
\end{tabular}
\end{center}

\begin{correctionbox}{Calcul du moment d'inertie d'un disque homogène}{}
Pour un disque de masse $M$, rayon $R$, densité surfacique $\sigma=M/(\pi R^2)$ :
\[
I_z=\iint r^2\dd m
=\int_0^R r^2\sigma(2\pi r\dd r)
=2\pi\sigma\int_0^R r^3\dd r.
\]
Donc
\[
I_z=2\pi\sigma\frac{R^4}{4}
=\frac{\pi\sigma R^4}{2}
=\frac12 MR^2.
\]
\end{correctionbox}

% ============================================================
\section{Chapitre 4 — Cinématique du solide indéformable}

\subsection{Champ des vitesses}

\begin{theoremBox}{Relation de Varignon des vitesses}{}
Pour un solide indéformable $\Solide$, il existe un vecteur rotation instantané $\vec\Omega$ tel que, pour tous points $A,B\in\Solide$ :
\[
\vec v_B=\vec v_A+\vec\Omega\wedge\vect{AB}.
\]
\end{theoremBox}

\begin{proof}
La distance $AB$ est constante :
\[
\norm{\vect{AB}}^2=\text{constante}.
\]
Donc
\[
\der{}{t}\left(\vect{AB}\cdot\vect{AB}\right)=2\vect{AB}\cdot(\vec v_B-\vec v_A)=0.
\]
Ainsi $\vec v_B-\vec v_A$ est orthogonal à $\vect{AB}$. On montre alors qu'il existe un vecteur $\vec\Omega$ indépendant du couple $(A,B)$ tel que
\[
\vec v_B-\vec v_A=\vec\Omega\wedge\vect{AB}.
\]
\end{proof}

\begin{center}
\begin{tikzpicture}[scale=1.1]
\draw[thick,fill=blue!7] (0,0) -- (3,0.5) -- (2.5,2) -- (-0.3,1.4) -- cycle;
\coordinate (A) at (0.7,0.7);
\coordinate (B) at (2.2,1.3);
\fill (A) circle (2pt) node[below]{$A$};
\fill (B) circle (2pt) node[above]{$B$};
\draw[->,very thick,red] (A)--++(1.0,0.2) node[right]{$\vec v_A$};
\draw[->,very thick,red] (B)--++(0.6,0.95) node[right]{$\vec v_B$};
\draw[->,very thick,vertprepa] (1.5,1)--++(0,1.2) node[above]{$\vec\Omega$};
\end{tikzpicture}
\end{center}

\subsection{Accélération dans un solide}

Si $A$ est un point du solide :
\[
\vec v_M=\vec v_A+\vec\Omega\wedge\vect{AM}.
\]
En dérivant dans un référentiel galiléen :
\[
\vec a_M
=
\vec a_A+\dot{\vec\Omega}\wedge\vect{AM}
+\vec\Omega\wedge(\vec\Omega\wedge\vect{AM}).
\]

\begin{warningbox}{Terme centripète}{}
Le terme
\[
\vec\Omega\wedge(\vec\Omega\wedge\vect{AM})
\]
ne disparaît pas lorsque $\vec\Omega$ est constant. Il correspond à l'accélération centripète.
\end{warningbox}

\subsection{Centre instantané de rotation dans le cas plan}

Dans un mouvement plan, s'il existe un point $I$ tel que $\vec v_I=\vec 0$, alors
\[
\vec v_M=\vec\Omega\wedge\vect{IM}.
\]
Le mouvement est instantanément équivalent à une rotation autour de $I$.

% ============================================================
\section{Chapitre 5 — Géométrie des masses et tenseur d'inertie}

\subsection{Tenseur d'inertie}

\begin{definitionbox}{Tenseur d'inertie en un point $O$}{}
Le tenseur d'inertie d'un solide par rapport au point $O$ est l'application linéaire définie par
\[
\vec L_O=\mat I_O\vec\Omega
\]
lorsque $O$ est fixe ou lorsque l'on calcule le moment cinétique propre autour de $G$.
Dans une base orthonormée :
\[
\mat I_O=
\begin{pmatrix}
A & -F & -E\\
-F & B & -D\\
-E & -D & C
\end{pmatrix},
\]
avec
\[
A=\int(y^2+z^2)\dd m,\quad
B=\int(x^2+z^2)\dd m,\quad
C=\int(x^2+y^2)\dd m,
\]
\[
D=\int yz\dd m,\quad
E=\int xz\dd m,\quad
F=\int xy\dd m.
\]
\end{definitionbox}

\begin{intuitionbox}{Sens du tenseur d'inertie}{}
Le moment d'inertie scalaire suffit pour une rotation autour d'un axe fixe connu.
En 3D, l'axe de rotation peut changer : la relation entre $\vec L$ et $\vec\Omega$ n'est plus forcément colinéaire.
Le tenseur d'inertie encode toute la géométrie de la masse.
\end{intuitionbox}

\begin{theoremBox}{Axes principaux d'inertie}{}
Comme $\mat I_O$ est une matrice symétrique réelle, elle est diagonalisable dans une base orthonormée.
Il existe donc une base d'axes principaux telle que
\[
\mat I_O=
\begin{pmatrix}
I_1&0&0\\
0&I_2&0\\
0&0&I_3
\end{pmatrix}.
\]
\end{theoremBox}

\subsection{Théorème de Huygens}

\begin{theoremBox}{Théorème de Huygens}{}
Soit $\Delta_G$ un axe passant par le centre de masse $G$ et $\Delta$ un axe parallèle à $\Delta_G$, situé à distance $d$.
Alors
\[
I_\Delta=I_{\Delta_G}+Md^2.
\]
\end{theoremBox}

\begin{proof}
Pour un point $M$ du solide :
\[
r_\Delta^2=r_{\Delta_G}^2+d^2+2d\,u,
\]
où $u$ est une coordonnée transverse relative à $G$.
En intégrant :
\[
I_\Delta=\int r_{\Delta_G}^2\dd m+Md^2+2d\int u\dd m.
\]
Le dernier terme est nul puisque $G$ est le centre de masse.
\end{proof}

\begin{center}
\begin{tikzpicture}[scale=1.1]
\draw[->,thick] (-0.5,0)--(4,0) node[right]{$x$};
\draw[->,thick] (0,-0.5)--(0,3) node[above]{$y$};
\draw[thick,fill=orange!12] (1.5,1.5) ellipse (1 and 0.65);
\fill (1.5,1.5) circle (2pt) node[below]{$G$};
\draw[very thick,blue] (1.5,0.3)--(1.5,2.8) node[above]{$\Delta_G$};
\draw[very thick,red] (2.6,0.3)--(2.6,2.8) node[above]{$\Delta$};
\draw[<->] (1.5,0.6)--(2.6,0.6) node[midway,below]{$d$};
\end{tikzpicture}
\end{center}

% ============================================================
\section{Chapitre 6 — Dynamique du solide indéformable}

\subsection{Torseur cinétique et torseur dynamique}

\begin{definitionbox}{Torseur cinétique}{}
Le torseur cinétique d'un solide $\Solide$ dans un référentiel $\Rep$ est
\[
\mathcal C_{\Solide/\Rep}
=
\left\{
\begin{array}{l}
\vec R_c=M\vec v_G,\\
\vec L_A=\displaystyle\int_{\Solide}\vect{AM}\wedge\vec v_M\,\dd m.
\end{array}
\right.
\]
\end{definitionbox}

\begin{definitionbox}{Torseur dynamique}{}
Le torseur dynamique est
\[
\mathcal D_{\Solide/\Rep}
=
\left\{
\begin{array}{l}
\vec R_d=M\vec a_G,\\
\vec \delta_A=\displaystyle\int_{\Solide}\vect{AM}\wedge\vec a_M\,\dd m.
\end{array}
\right.
\]
\end{definitionbox}

\begin{theoremBox}{Principe fondamental de la dynamique du solide}{}
Dans un référentiel galiléen, pour un solide $\Solide$ :
\[
M\vec a_G=\sum \vec F_{\rm ext},
\]
et, pour un point fixe $A$ du référentiel galiléen :
\[
\vec \delta_A=\sum \vec M_A^{\rm ext}.
\]
\end{theoremBox}

\subsection{Rotation autour d'un axe fixe}

Si le solide tourne autour d'un axe fixe $\Delta$ avec vitesse angulaire $\omega$ :
\[
\sum M_\Delta^{\rm ext}=I_\Delta\dot\omega.
\]

\begin{examplebox}{Pendule pesant}{}
Un solide de masse $M$, de centre de masse $G$, est mobile autour d'un axe horizontal passant par $O$.
On note $a=OG$, $I_O$ le moment d'inertie autour de l'axe, et $\theta$ l'angle avec la verticale.
Le moment du poids vaut
\[
M_O(P)=-Mga\sin\theta.
\]
L'équation du mouvement est
\[
I_O\ddot\theta+Mga\sin\theta=0.
\]
Pour petits angles :
\[
\ddot\theta+\frac{Mga}{I_O}\theta=0.
\]
Donc
\[
T=2\pi\sqrt{\frac{I_O}{Mga}}.
\]
\end{examplebox}

\subsection{Équations d'Euler}

Dans une base principale liée au solide :
\[
\vec L_G=I_1\Omega_1\vec e_1+I_2\Omega_2\vec e_2+I_3\Omega_3\vec e_3.
\]

La dérivation dans une base mobile donne :
\[
\left(\der{\vec L}{t}\right)_{\Rep}
=
\left(\der{\vec L}{t}\right)_{\text{solide}}
+
\vec\Omega\wedge\vec L.
\]

\begin{theoremBox}{Équations d'Euler}{}
Dans une base principale liée au solide :
\[
\begin{cases}
I_1\dot\Omega_1+(I_3-I_2)\Omega_2\Omega_3=M_1,\\
I_2\dot\Omega_2+(I_1-I_3)\Omega_3\Omega_1=M_2,\\
I_3\dot\Omega_3+(I_2-I_1)\Omega_1\Omega_2=M_3.
\end{cases}
\]
\end{theoremBox}

% ============================================================
\section{Chapitre 7 — Énergie d'un solide}

\begin{theoremBox}{Énergie cinétique d'un solide}{}
Pour un solide indéformable :
\[
E_c=\frac12 Mv_G^2+\frac12\vec\Omega\cdot\mat I_G\vec\Omega.
\]
\end{theoremBox}

\begin{proof}
On écrit
\[
\vec v_M=\vec v_G+\vec\Omega\wedge\vect{GM}.
\]
Alors
\[
E_c=\frac12\int v_M^2\dd m.
\]
En développant :
\[
E_c=\frac12 Mv_G^2
+\vec v_G\cdot\int(\vec\Omega\wedge\vect{GM})\dd m
+\frac12\int\norm{\vec\Omega\wedge\vect{GM}}^2\dd m.
\]
Le terme croisé est nul car
\[
\int\vect{GM}\dd m=\vec 0.
\]
Le dernier terme vaut
\[
\frac12\vec\Omega\cdot\mat I_G\vec\Omega.
\]
\end{proof}

\begin{examplebox}{Disque roulant sans glissement}{}
Un disque plein de masse $M$ et rayon $R$ roule sans glisser à vitesse $v_G$.
Comme
\[
I_G=\frac12MR^2,
\qquad
v_G=R\omega,
\]
on obtient
\[
E_c=\frac12Mv_G^2+\frac12I_G\omega^2
=\frac12Mv_G^2+\frac14Mv_G^2
=\frac34Mv_G^2.
\]
\end{examplebox}

% ============================================================
\section{Chapitre 8 — Liaisons mécaniques et modélisation}

\subsection{Liaisons usuelles}

\begin{center}
\begin{tabular}{lll}
\toprule
Liaison & Degrés de liberté bloqués & Exemple \\
\midrule
Pivot & 5 & Porte, roue sur axe \\
Glissière & 5 & Tiroir, vérin \\
Rotule & 3 & Articulation sphérique \\
Appui plan & 3 & Bloc posé sur table \\
Ponctuelle & 1 & Contact bille-plan idéal \\
\bottomrule
\end{tabular}
\end{center}

\begin{definitionbox}{Liaison parfaite}{}
Une liaison est dite parfaite si les actions de liaison ne dissipent pas d'énergie.
Autrement dit, leur puissance virtuelle est nulle pour les mouvements compatibles avec la liaison.
\end{definitionbox}

\begin{methodbox}{Isoler un solide}{}
Pour étudier un solide :
\begin{enumerate}
  \item dessiner le solide seul ;
  \item représenter toutes les actions extérieures ;
  \item identifier les inconnues de liaison ;
  \item choisir le point de calcul des moments ;
  \item écrire les équations de la dynamique ou de la statique ;
  \item vérifier le nombre d'équations et d'inconnues.
\end{enumerate}
\end{methodbox}

% ============================================================
\section{Chapitre 9 — Roulement sans glissement}

\begin{definitionbox}{Condition de roulement sans glissement}{}
Un solide roule sans glisser sur un support fixe si le point de contact $I$ a une vitesse nulle dans le référentiel du support :
\[
\vec v_I=\vec 0.
\]
Pour une roue de rayon $R$ :
\[
v_G=R\omega.
\]
\end{definitionbox}

\begin{center}
\begin{tikzpicture}[scale=1.1]
\draw[thick] (-0.5,0)--(5,0);
\draw[fill=blue!7,thick] (2,1) circle (1);
\fill (2,1) circle (2pt) node[above]{$G$};
\fill (2,0) circle (2pt) node[below]{$I$};
\draw[->,very thick,red] (2,1)--(3.2,1) node[right]{$\vec v_G$};
\draw[->,very thick,vertprepa] (2,1)--++(0,1.1) node[above]{$\vec\Omega$};
\draw[<->] (2,0)--(2,1) node[midway,right]{$R$};
\node at (2,-0.45) {$\vec v_I=\vec 0$};
\end{tikzpicture}
\end{center}

\begin{examplebox}{Cylindre descendant un plan incliné}{}
Un cylindre de masse $M$, rayon $R$, moment $I_G=\alpha MR^2$, descend un plan incliné d'angle $\theta$ sans glisser.

L'énergie donne
\[
Mg h=\frac12Mv^2+\frac12I_G\omega^2.
\]
Avec $\omega=v/R$ :
\[
Mg h=\frac12Mv^2(1+\alpha).
\]
Donc
\[
v=\sqrt{\frac{2gh}{1+\alpha}}.
\]
Pour un cylindre plein $\alpha=1/2$ :
\[
v=\sqrt{\frac{4gh}{3}}.
\]
Pour un cylindre creux mince $\alpha=1$ :
\[
v=\sqrt{gh}.
\]
Le cylindre plein descend plus vite.
\end{examplebox}

% ============================================================
\section{Chapitre 10 — Petits mouvements et stabilité}

\begin{definitionbox}{Équilibre stable}{}
Une position d'équilibre est stable si une petite perturbation produit une force ou un moment de rappel.
Énergétiquement, un équilibre est stable lorsque l'énergie potentielle possède un minimum local.
\end{definitionbox}

\begin{methodbox}{Étude des petits mouvements}{}
\begin{enumerate}
  \item Trouver la position d'équilibre.
  \item Écrire l'équation exacte du mouvement.
  \item Linéariser autour de l'équilibre.
  \item Obtenir une équation de type
  \[
  \ddot x+\omega_0^2x=0.
  \]
  \item En déduire la période :
  \[
  T=\frac{2\pi}{\omega_0}.
  \]
\end{enumerate}
\end{methodbox}

\begin{examplebox}{Pendule physique}{}
Pour le pendule pesant :
\[
I_O\ddot\theta+Mga\sin\theta=0.
\]
Pour $\theta$ petit :
\[
\sin\theta\simeq \theta,
\]
donc
\[
\ddot\theta+\frac{Mga}{I_O}\theta=0.
\]
L'équilibre bas est stable.
\end{examplebox}

% ============================================================
\section{Chapitre 11 — Problèmes contextualisés avancés}

% ------------------------------------------------------------
\subsection{Problème 1 — Jeu de bar : cylindre rempli d'un liquide visqueux}

\begin{exercisebox}{Énoncé}{}
Un cylindre vertical transparent de rayon intérieur $R_c=\SI{8.0}{cm}$ et de hauteur utile $H=\SI{60}{cm}$ est rempli d'un liquide de masse volumique
\[
\rho_f=\SI{1050}{kg.m^{-3}}
\]
et de viscosité dynamique
\[
\eta=\SI{0.80}{Pa.s}.
\]
On lâche une bille sphérique de rayon
\[
a=\SI{5.0}{mm}
\]
et de masse volumique
\[
\rho_b=\SI{7800}{kg.m^{-3}}.
\]
Une cible circulaire de diamètre \SI{3.0}{cm} est placée au fond, à une distance horizontale $x_c=\SI{4.0}{cm}$ de la verticale de lâcher.

On suppose d'abord que la traînée est linéaire :
\[
\vec f=-6\pi\eta a\vec v.
\]

\begin{enumerate}
  \item Établir l'équation différentielle du mouvement vertical.
  \item Résoudre $v_z(t)$ en supposant $v_z(0)=0$.
  \item Déterminer la vitesse limite.
  \item Estimer le temps de chute.
  \item Étudier le mouvement horizontal si la bille reçoit une vitesse initiale $v_{x0}$.
  \item Déterminer $v_{x0}$ pour atteindre la cible.
  \item Étudier la sensibilité à une erreur de vitesse initiale.
  \item Conclure sur une stratégie optimale neutre physiquement.
\end{enumerate}
\end{exercisebox}

\begin{center}
\begin{tikzpicture}[scale=1.1]
\draw[thick] (0,0) rectangle (3,6);
\fill[blue!10] (0,0) rectangle (3,6);
\draw[thick,fill=gray!20] (1.5,5.4) circle (0.12);
\draw[->,red,thick] (1.5,5.4)--(2.0,5.4) node[right]{$v_{x0}$};
\draw[->,vertprepa,thick] (1.5,5.4)--(1.5,4.8) node[right]{$\vec P-\vec\Pi$};
\draw[fill=orange!50] (2.3,0.15) circle (0.18);
\node at (2.3,-0.25) {cible};
\draw[dashed] (1.5,5.4)--(1.5,0);
\draw[<->] (3.25,0)--(3.25,6) node[midway,right]{$H$};
\draw[<->] (1.5,0.45)--(2.3,0.45) node[midway,above]{$x_c$};
\end{tikzpicture}
\end{center}

\begin{correctionbox}{Résolution détaillée}{}
On oriente $z$ vers le bas. Le volume de la bille est
\[
V=\frac43\pi a^3.
\]
Sa masse vaut
\[
m=\rho_b V.
\]
Les forces verticales sont :
\[
P=mg,\qquad
\Pi=\rho_f Vg,\qquad
f_z=-6\pi\eta a v_z.
\]
L'équation est donc
\[
m\der{v_z}{t}=mg-\rho_fVg-6\pi\eta a v_z.
\]
Comme $m=\rho_bV$ :
\[
\der{v_z}{t}+\frac{1}{\tau}v_z
=
g\left(1-\frac{\rho_f}{\rho_b}\right),
\]
avec
\[
\tau=\frac{m}{6\pi\eta a}
=\frac{2\rho_ba^2}{9\eta}.
\]
La solution avec $v_z(0)=0$ est
\[
v_z(t)=v_\infty(1-e^{-t/\tau}),
\]
où
\[
v_\infty
=
g\left(1-\frac{\rho_f}{\rho_b}\right)\tau
=
\frac{2a^2g(\rho_b-\rho_f)}{9\eta}.
\]

Application numérique :
\[
v_\infty
=
\frac{2(5.0\times10^{-3})^2\times9.81(7800-1050)}
{9\times0.80}
\simeq \SI{0.46}{m.s^{-1}}.
\]
Le temps caractéristique vaut
\[
\tau=\frac{2\times7800\times(5.0\times10^{-3})^2}{9\times0.80}
\simeq \SI{0.054}{s}.
\]
Comme $\tau$ est très petit devant le temps de chute, on approxime
\[
t_f\simeq \frac{H}{v_\infty}
=\frac{0.60}{0.46}
\simeq \SI{1.30}{s}.
\]

Horizontalement :
\[
m\der{v_x}{t}=-6\pi\eta a v_x,
\]
donc
\[
v_x(t)=v_{x0}e^{-t/\tau},
\]
et
\[
x(t)=v_{x0}\tau(1-e^{-t/\tau}).
\]
Comme $t_f\gg \tau$ :
\[
x_f\simeq v_{x0}\tau.
\]
Pour viser $x_c=\SI{4.0}{cm}$ :
\[
v_{x0}\simeq \frac{x_c}{\tau}
=\frac{0.040}{0.054}
\simeq \SI{0.74}{m.s^{-1}}.
\]

La sensibilité est
\[
\delta x_f\simeq \tau\,\delta v_{x0}.
\]
Pour rester dans une cible de rayon \SI{1.5}{cm} :
\[
|\delta v_{x0}|\lesssim \frac{0.015}{0.054}
\simeq \SI{0.28}{m.s^{-1}}.
\]

Conclusion physique neutre :
\begin{itemize}
  \item dans un fluide très visqueux, la vitesse horizontale est rapidement amortie ;
  \item la bille doit recevoir une vitesse initiale horizontale assez précise ;
  \item une bille plus dense ou plus grande tombe plus vite, mais la traînée change aussi ;
  \item la meilleure stratégie consiste à calibrer expérimentalement le couple bille-liquide-cible, puis à reproduire une vitesse initiale faible mais contrôlée.
\end{itemize}
\end{correctionbox}

\begin{center}
\begin{tikzpicture}
\begin{axis}[
width=10cm,height=6cm,
xlabel={$t$ (s)},
ylabel={$v_z(t)$ (m/s)},
grid=both,
domain=0:0.35,
samples=100]
\addplot {0.46*(1-exp(-x/0.054))};
\end{axis}
\end{tikzpicture}
\captionof{figure}{Approche exponentielle de la vitesse limite dans un fluide visqueux.}
\end{center}

% ------------------------------------------------------------
\subsection{Problème 2 — Cylindre roulant sur un plan incliné}

\begin{exercisebox}{Énoncé}{}
Un cylindre de masse $M=\SI{2.0}{kg}$ et de rayon $R=\SI{10}{cm}$ descend sans glisser un plan incliné d'angle $\theta=25^\circ$ et de longueur $L=\SI{2.0}{m}$.
Comparer un cylindre plein et un cylindre creux mince.
\end{exercisebox}

\begin{center}
\begin{tikzpicture}[scale=1.1]
\draw[thick] (0,0)--(5,2);
\draw[thick,fill=blue!8,rotate around={21.8:(2.2,1.0)}] (2.2,1.0) circle (0.45);
\draw[->,red,thick] (2.2,1.0)--++(0.9,0.36) node[right]{$\vec v_G$};
\draw[->,thick] (0.8,0.1) arc[start angle=0,end angle=21.8,radius=0.8];
\node at (1.25,0.28){$\theta$};
\end{tikzpicture}
\end{center}

\begin{correctionbox}{Correction}{}
La perte d'énergie potentielle est
\[
Mgh=MgL\sin\theta.
\]
L'énergie cinétique en bas est
\[
E_c=\frac12Mv^2+\frac12I_G\omega^2
=\frac12Mv^2(1+\alpha),
\]
avec $I_G=\alpha MR^2$ et $\omega=v/R$.
Donc
\[
v=\sqrt{\frac{2gL\sin\theta}{1+\alpha}}.
\]
Pour le cylindre plein, $\alpha=1/2$ :
\[
v_{\rm plein}
=
\sqrt{\frac{2gL\sin\theta}{3/2}}
\simeq \SI{3.32}{m.s^{-1}}.
\]
Pour le cylindre creux mince, $\alpha=1$ :
\[
v_{\rm creux}
=
\sqrt{gL\sin\theta}
\simeq \SI{2.88}{m.s^{-1}}.
\]
Le cylindre plein arrive en premier car il possède un moment d'inertie plus faible relativement à sa masse.
\end{correctionbox}

% ------------------------------------------------------------
\subsection{Problème 3 — Porte qui claque}

\begin{exercisebox}{Énoncé}{}
Une porte rectangulaire de masse $M=\SI{25}{kg}$, largeur $L=\SI{0.90}{m}$, hauteur \SI{2.1}{m}, tourne autour d'un axe vertical.
Un coup de vent exerce une force moyenne $F=\SI{20}{N}$ au centre de la porte, perpendiculairement au plan de la porte.
On néglige les frottements pendant \SI{0.50}{s}. Estimer la vitesse angulaire atteinte.
\end{exercisebox}

\begin{center}
\begin{tikzpicture}[scale=1.1]
\draw[very thick] (0,0)--(0,3);
\draw[fill=orange!15,thick] (0,0.3) rectangle (2.2,2.7);
\draw[->,red,very thick] (1.1,1.5)--(1.1,2.4) node[above]{$\vec F$};
\draw[<->] (0,0.1)--(2.2,0.1) node[midway,below]{$L$};
\node[left] at (0,1.5){axe};
\end{tikzpicture}
\end{center}

\begin{correctionbox}{Correction}{}
Le moment de la force par rapport à l'axe vaut approximativement
\[
M_\Delta=F\frac L2.
\]
Le moment d'inertie d'une porte autour d'un bord vertical est
\[
I_\Delta=\frac13ML^2.
\]
Donc
\[
I_\Delta\dot\omega=F\frac L2.
\]
Ainsi
\[
\dot\omega=\frac{FL/2}{ML^2/3}
=\frac{3F}{2ML}.
\]
Numériquement :
\[
\dot\omega=\frac{3\times20}{2\times25\times0.90}
\simeq \SI{1.33}{rad.s^{-2}}.
\]
Après \SI{0.50}{s} :
\[
\omega\simeq \SI{0.67}{rad.s^{-1}}.
\]
\end{correctionbox}

% ------------------------------------------------------------
\subsection{Problème 4 — Skateboard ou trottinette}

\begin{exercisebox}{Énoncé}{}
Un rider de masse totale $M=\SI{75}{kg}$ avance à \SI{5.0}{m.s^{-1}} sur une trottinette.
Chaque roue a masse $m=\SI{0.30}{kg}$, rayon $R=\SI{6.0}{cm}$ et moment d'inertie $I=\frac12mR^2$.
Calculer l'énergie cinétique totale en tenant compte de la rotation des deux roues.
\end{exercisebox}

\begin{correctionbox}{Correction}{}
L'énergie de translation du système est
\[
E_{\rm trans}=\frac12Mv^2
=
\frac12\times75\times5^2
=\SI{937.5}{J}.
\]
Pour une roue :
\[
E_{\rm rot,1}=\frac12I\omega^2
=\frac12\cdot\frac12mR^2\cdot\frac{v^2}{R^2}
=\frac14mv^2.
\]
Pour deux roues :
\[
E_{\rm rot}=2\cdot\frac14mv^2
=\frac12mv^2
=\frac12\times0.30\times25
=\SI{3.75}{J}.
\]
La rotation des petites roues est négligeable devant la translation du rider, mais elle devient importante si les roues sont lourdes.
\end{correctionbox}

% ------------------------------------------------------------
\subsection{Problème 5 — Verre qui bascule}

\begin{exercisebox}{Énoncé}{}
On modélise un verre par un cylindre de rayon $R=\SI{3.0}{cm}$ et de hauteur $H=\SI{12}{cm}$.
Son centre de masse est à hauteur $h_G=\SI{5.0}{cm}$.
Déterminer l'angle critique de basculement.
\end{exercisebox}

\begin{center}
\begin{tikzpicture}[scale=1.1]
\draw[thick,rotate around={18:(0,0)}] (-0.6,0) rectangle (0.6,2.4);
\fill[red,rotate around={18:(0,0)}] (0,1.0) circle (2pt) node[right]{$G$};
\draw[dashed] (0,0)--(0,2.6);
\draw[->,red,thick,rotate around={18:(0,0)}] (0,1.0)--(0,0.2) node[right]{$\vec P$};
\draw[->] (1.3,0) arc[start angle=0,end angle=18,radius=1.3];
\node at (1.5,0.25){$\alpha$};
\end{tikzpicture}
\end{center}

\begin{correctionbox}{Correction}{}
Le basculement commence lorsque la verticale du centre de masse passe par le bord de la base.
Ainsi :
\[
\tan\alpha_c=\frac{R}{h_G}.
\]
Donc
\[
\alpha_c=\arctan\left(\frac{0.030}{0.050}\right)
\simeq 31^\circ.
\]
Si le verre est rempli, $h_G$ change. Un centre de masse plus bas augmente la stabilité.
\end{correctionbox}

% ------------------------------------------------------------
\subsection{Problème 6 — Hand spinner}

\begin{exercisebox}{Énoncé}{}
Un hand spinner de moment d'inertie $I=\SI{1.8e-4}{kg.m^2}$ tourne initialement à $\omega_0=\SI{80}{rad.s^{-1}}$.
Le couple de frottement est supposé constant :
\[
C_f=\SI{2.0e-4}{N.m}.
\]
Déterminer le temps d'arrêt.
\end{exercisebox}

\begin{correctionbox}{Correction}{}
L'équation dynamique est
\[
I\dot\omega=-C_f.
\]
Donc
\[
\omega(t)=\omega_0-\frac{C_f}{I}t.
\]
Le temps d'arrêt vérifie $\omega(t_a)=0$ :
\[
t_a=\frac{I\omega_0}{C_f}
=
\frac{1.8\times10^{-4}\times80}{2.0\times10^{-4}}
=\SI{72}{s}.
\]
\end{correctionbox}

% ------------------------------------------------------------
\subsection{Problème 7 — Échelle contre un mur}

\begin{exercisebox}{Énoncé}{}
Une échelle homogène de longueur $L$ et masse $M$ repose contre un mur vertical sans frottement.
Le sol possède un coefficient de frottement statique $\mu_s$.
Déterminer l'angle minimal $\theta$ avec le sol pour que l'échelle ne glisse pas.
\end{exercisebox}

\begin{center}
\begin{tikzpicture}[scale=1.1]
\draw[thick] (0,0)--(5,0);
\draw[thick] (5,0)--(5,4);
\draw[very thick,brown] (1,0)--(5,3);
\draw[->,red,thick] (3,1.5)--(3,0.5) node[right]{$Mg$};
\draw[->,blue,thick] (1,0)--(1,0.8) node[left]{$N_s$};
\draw[->,blue,thick] (1,0)--(1.0+0.8,0) node[below]{$T_s$};
\draw[->,blue,thick] (5,3)--(4.2,3) node[above]{$N_m$};
\draw[->] (1.8,0) arc[start angle=0,end angle=36.9,radius=0.8];
\node at (2.1,0.25){$\theta$};
\end{tikzpicture}
\end{center}

\begin{correctionbox}{Correction}{}
Les forces sont :
\[
Mg,\quad N_s,\quad T_s,\quad N_m.
\]
Équilibre horizontal :
\[
T_s=N_m.
\]
Équilibre vertical :
\[
N_s=Mg.
\]
Moment au point de contact avec le sol :
\[
N_m L\sin\theta=Mg\frac L2\cos\theta.
\]
Donc
\[
N_m=\frac{Mg}{2}\cot\theta.
\]
La condition de non-glissement est
\[
T_s\leq \mu_s N_s.
\]
Ainsi
\[
\frac{Mg}{2}\cot\theta\leq \mu_s Mg,
\]
d'où
\[
\tan\theta\geq \frac{1}{2\mu_s}.
\]
Finalement :
\[
\theta_{\min}=\arctan\left(\frac{1}{2\mu_s}\right).
\]
\end{correctionbox}

% ------------------------------------------------------------
\subsection{Problème 8 — Solide composite}

\begin{exercisebox}{Énoncé}{}
Un solide composite est constitué de deux disques homogènes coaxiaux soudés :
\[
M_1=\SI{2.0}{kg},\quad R_1=\SI{10}{cm},
\qquad
M_2=\SI{1.0}{kg},\quad R_2=\SI{20}{cm}.
\]
Calculer le moment d'inertie autour de l'axe commun.
\end{exercisebox}

\begin{correctionbox}{Correction}{}
Pour chaque disque :
\[
I_i=\frac12M_iR_i^2.
\]
Donc
\[
I=I_1+I_2
=
\frac12M_1R_1^2+\frac12M_2R_2^2.
\]
Numériquement :
\[
I=\frac12\times2.0\times0.10^2
+\frac12\times1.0\times0.20^2
=\SI{0.010}{kg.m^2}+\SI{0.020}{kg.m^2}
=\SI{0.030}{kg.m^2}.
\]
\end{correctionbox}

% ============================================================
\section{Banque d'exercices progressifs}

\subsection*{Chapitre 0 — Rappels}
\begin{enumerate}
  \item \niveau{$\star$} Calculer l'accélération d'un point soumis à une force constante.
  \item \niveau{$\star$} Établir l'équation horaire d'une chute libre sans frottement.
  \item \niveau{$\star\star$} Étudier un ressort vertical avec gravité.
  \item \niveau{$\star\star$} Calculer la vitesse minimale pour franchir un looping.
  \item \niveau{$\star\star\star$} Étudier un projectile avec frottement linéaire.
\end{enumerate}

\subsection*{Chapitre 1 — Point matériel}
\begin{enumerate}
  \item \niveau{$\star$} Retrouver les expressions de $\vec v$ et $\vec a$ en coordonnées polaires.
  \item \niveau{$\star\star$} Montrer que le moment cinétique est conservé pour une force centrale.
  \item \niveau{$\star\star\star$} Étudier une orbite circulaire sous force gravitationnelle.
  \item \niveau{$\star\star\star\star$} Étudier la stabilité d'une orbite circulaire.
\end{enumerate}

\subsection*{Chapitre 2 — Systèmes discrets}
\begin{enumerate}
  \item \niveau{$\star$} Calculer le centre de masse de deux masses ponctuelles.
  \item \niveau{$\star\star$} Démontrer le théorème de la résultante cinétique.
  \item \niveau{$\star\star\star$} Étudier une collision parfaitement inélastique.
  \item \niveau{$\star\star\star\star$} Étudier une explosion interne avec conservation de la quantité de mouvement.
\end{enumerate}

\subsection*{Chapitre 3 — Continu}
\begin{enumerate}
  \item \niveau{$\star$} Calculer le centre de masse d'une tige homogène.
  \item \niveau{$\star\star$} Calculer le centre de masse d'une plaque triangulaire.
  \item \niveau{$\star\star\star$} Calculer le moment d'inertie d'un disque.
  \item \niveau{$\star\star\star\star$} Calculer le moment d'inertie d'un cône homogène.
\end{enumerate}

\subsection*{Chapitre 4 — Cinématique du solide}
\begin{enumerate}
  \item \niveau{$\star$} Donner le champ des vitesses d'une barre en rotation.
  \item \niveau{$\star\star$} Déterminer le centre instantané de rotation d'une roue.
  \item \niveau{$\star\star\star$} Calculer l'accélération d'un point d'une roue.
  \item \niveau{$\star\star\star\star$} Étudier une barre glissant entre deux murs perpendiculaires.
\end{enumerate}

\subsection*{Chapitre 5 — Inertie}
\begin{enumerate}
  \item \niveau{$\star$} Utiliser Huygens pour une tige autour d'une extrémité.
  \item \niveau{$\star\star$} Calculer le tenseur d'inertie d'un parallélépipède rectangle.
  \item \niveau{$\star\star\star$} Trouver les axes principaux d'un solide plan symétrique.
  \item \niveau{$\star\star\star\star$} Diagonaliser un tenseur d'inertie non diagonal.
\end{enumerate}

\subsection*{Chapitre 6 — Dynamique}
\begin{enumerate}
  \item \niveau{$\star$} Étudier une porte soumise à un couple constant.
  \item \niveau{$\star\star$} Établir l'équation d'un pendule pesant.
  \item \niveau{$\star\star\star$} Étudier un disque entraîné par un fil.
  \item \niveau{$\star\star\star\star$} Exploiter les équations d'Euler pour une toupie symétrique.
\end{enumerate}

\subsection*{Chapitre 7 — Énergie}
\begin{enumerate}
  \item \niveau{$\star$} Calculer l'énergie cinétique d'un disque roulant.
  \item \niveau{$\star\star$} Comparer disque, anneau et sphère roulant sans glisser.
  \item \niveau{$\star\star\star$} Étudier une bille roulant dans une cuvette.
\end{enumerate}

\subsection*{Chapitre 8 — Liaisons}
\begin{enumerate}
  \item \niveau{$\star$} Identifier les inconnues d'une liaison pivot.
  \item \niveau{$\star\star$} Faire le bilan d'actions sur une échelle.
  \item \niveau{$\star\star\star$} Étudier un mécanisme bielle-manivelle simplifié.
\end{enumerate}

\subsection*{Chapitre 9 — Roulement}
\begin{enumerate}
  \item \niveau{$\star$} Écrire la condition $v_G=R\omega$.
  \item \niveau{$\star\star$} Calculer l'accélération d'un cylindre sur plan incliné.
  \item \niveau{$\star\star\star$} Déterminer la force de frottement statique nécessaire.
  \item \niveau{$\star\star\star\star$} Étudier la transition glissement-roulement.
\end{enumerate}

\subsection*{Chapitre 10 — Stabilité}
\begin{enumerate}
  \item \niveau{$\star$} Linéariser $\sin\theta$.
  \item \niveau{$\star\star$} Étudier la stabilité d'un verre incliné.
  \item \niveau{$\star\star\star$} Trouver la période d'un pendule physique.
  \item \niveau{$\star\star\star\star$} Étudier des petites oscillations couplées de deux solides.
\end{enumerate}

% ============================================================
\section{Formulaire final}

\begin{multicols}{2}

\subsection*{Point matériel}
\[
\vec v=\der{\vec r}{t},\qquad
\vec a=\der{\vec v}{t}
\]
\[
m\vec a=\sum\vec F
\]
\[
\vec L_O=\vect{OM}\wedge m\vec v
\]
\[
\der{\vec L_O}{t}=\sum \vec M_O
\]

\subsection*{Centre de masse}
\[
M=\sum_i m_i
\]
\[
\vect{OG}=\frac1M\sum_i m_i\vect{OM_i}
\]
\[
\vect{OG}=\frac1M\int \vect{OM}\dd m
\]

\subsection*{Solide}
\[
\vec v_B=\vec v_A+\vec\Omega\wedge\vect{AB}
\]
\[
\vec a_B=\vec a_A+\dot{\vec\Omega}\wedge\vect{AB}
+\vec\Omega\wedge(\vec\Omega\wedge\vect{AB})
\]

\subsection*{Inertie}
\[
I_\Delta=\int r_\perp^2\dd m
\]
\[
I_\Delta=I_{\Delta_G}+Md^2
\]
\[
\vec L_G=\mat I_G\vec\Omega
\]

\subsection*{Énergie}
\[
E_c=\frac12Mv_G^2+\frac12\vec\Omega\cdot\mat I_G\vec\Omega
\]

\subsection*{Roulement}
\[
\vec v_I=\vec 0
\]
\[
v_G=R\omega
\]

\subsection*{Euler}
\[
I_1\dot\Omega_1+(I_3-I_2)\Omega_2\Omega_3=M_1
\]
\[
I_2\dot\Omega_2+(I_1-I_3)\Omega_3\Omega_1=M_2
\]
\[
I_3\dot\Omega_3+(I_2-I_1)\Omega_1\Omega_2=M_3
\]

\end{multicols}

% ============================================================
\section*{Conclusion pédagogique}
\addcontentsline{toc}{section}{Conclusion pédagogique}

La mécanique du solide prolonge la mécanique du point en ajoutant une idée fondamentale :
la masse n'est plus concentrée en un point, mais répartie dans l'espace.
Cette répartition impose d'introduire le centre de masse, les moments d'inertie, le tenseur d'inertie,
les torseurs et la cinématique des rotations.

Le passage essentiel est donc :
\[
\text{point matériel}
\quad\longrightarrow\quad
\text{système discret}
\quad\longrightarrow\quad
\text{milieu continu}
\quad\longrightarrow\quad
\text{solide indéformable 3D}.
\]

Dans les problèmes réels, la difficulté n'est pas seulement de connaître les formules :
elle consiste à choisir le bon modèle, identifier les hypothèses, faire le bon bilan d'actions,
et interpréter les ordres de grandeur.

\end{document}